% Generated draft: requires mathematical and source review.
\documentclass[11pt]{article}
\usepackage[margin=25mm]{geometry}
\usepackage{fontspec}
\setmainfont{Latin Modern Roman}
\usepackage{amsmath,amssymb,mathtools,unicode-math}
\setmathfont{Latin Modern Math}
\usepackage{enumitem,needspace,xurl,hyperref}
\setlength{\parindent}{0pt}
\setlength{\parskip}{5pt}
\setlength{\emergencystretch}{3em}
\newcommand{\N}{\mathbb{N}}
\newcommand{\Z}{\mathbb{Z}}
\newcommand{\Q}{\mathbb{Q}}
\newcommand{\R}{\mathbb{R}}
\newcommand{\C}{\mathbb{C}}
\newcommand{\F}{\mathbb{F}}
\newcommand{\A}{\mathbb{A}}
\newcommand{\PP}{\mathbb{P}}
\newcommand{\E}{\mathbb{E}}
\newcommand{\eps}{\varepsilon}
\newcommand{\dd}{\,\mathrm{d}}
\begin{document}
\section*{UnsolvedMath ID: 701032}
\textbf{Source problem code:} catalogue entry
\subsection{Definitions and mathematical statement}
A finite group \(G\) is called soluble if its derived series, defined by \(G^{(0)}=G\) and \(G^{(n+1)}=[G^{(n)},G^{(n)}]\), reaches the trivial group. A soluble Fitting class is a nonempty class \(\mathcal F\) of finite soluble groups, closed under isomorphisms, satisfying:
\begin{enumerate}
\item If \(G\in\mathcal F\) and \(N\trianglelefteq G\), then \(N\in\mathcal F\).
\item If \(N_1,N_2\trianglelefteq G\) for a finite group \(G\), and \(N_1,N_2\in\mathcal F\), then \(N_1N_2\in\mathcal F\).
\end{enumerate}
Let \(\mathfrak L\) be the lattice of all soluble Fitting classes, ordered by inclusion. Its meet is \(\mathcal X\wedge\mathcal Y=\mathcal X\cap\mathcal Y\), and its join \(\mathcal X\vee\mathcal Y\) is the smallest soluble Fitting class containing \(\mathcal X\cup\mathcal Y\).

The question is whether \(\mathfrak L\) is modular: does
\[
\mathcal X\vee(\mathcal Y\wedge\mathcal Z)
=(\mathcal X\vee\mathcal Y)\wedge\mathcal Z
\]
hold for every \(\mathcal X,\mathcal Y,\mathcal Z\in\mathfrak L\) with \(\mathcal X\subseteq\mathcal Z\)? Here the standard nonempty convention for Fitting classes is used; the supplied statement does not explicitly give a definition. The supplied registry labels the problem open, but that status has not been independently verified.
\subsection{Short English statement}
Is the lattice, ordered by inclusion, of all Fitting classes consisting of finite soluble groups modular? Equivalently, must the displayed modular identity hold for every three such classes whenever \(\mathcal X\subseteq\mathcal Z\)?
\subsection{Sources}
\begin{itemize}
\item Supplied local UnsolvedMath registry record: permanent ID \(701032\), source problem code \texttt{catalogue entry}, title \emph{Kourovka Notebook Problem 14.47}. The record labels the problem \texttt{open} and was updated on \texttt{2026-10-16}; this is supplied registry metadata, not independent verification of its mathematical status.
\item Supplied attribution: proposed by S. F. Kamornikov and A. N. Skiba in 1999.
\item Supplied source: \emph{The Kourovka Notebook (Unsolved Problems in Group Theory)}, 14th issue (1999), Problem 14.47; the record also identifies the 21st issue edition (2026), arXiv:1401.0300v48. The collection description names E. I. Khukhro and V. D. Mazurov as editors. These bibliographic details are supplied attribution and have not been independently checked.
\end{itemize}
\end{document}
